\chapter{Discussion} \label{chapter:ch6_discussion}

The introduction asked whether frequently changing collaborative-robot tasks can be programmed more directly by specifying spatial intent in the real workspace instead of relying on hand-guiding or editing motion commands and tool actions in a teach pendant or offline programming workflow. Within the scope of this prototype, the result is positive for selected task-oriented workflows.

This chapter interprets that result rather than recapping the implementation. The thesis does not propose a general robot programming environment. It shows that a selected use case can combine workspace sensing, spatial input, task-specific interpretation, and robot execution into one authoring workflow, and it also makes clear where that result currently stops.

\section{Prototype Result}

The main result is an end-to-end prototype for in-workspace spatial authoring on a collaborative robot. The implemented system combines a tracked 6DoF input device, stereo-vision-based scene sensing, task-specific interpretation, and robot execution behind a shared modular runtime with explicit interfaces. In practical terms, the prototype covers the full path from setup and calibration through authoring, confirmation, visualization, execution, and persistence.

At prototype scope, this closes the three goals stated in the introduction. For \goalref{goal:1}, the thesis demonstrates that poses and trajectories can be authored directly in the workspace and turned into executable robot behavior through task-specific interpretation. For \goalref{goal:2}, it provides a modular system in which different use cases reuse the same sensing, orchestration, and execution infrastructure. For \goalref{goal:3}, it demonstrates that approach on a real robot with two representative task families, pick-and-place and seam welding, while keeping the evaluation qualitative.

The important result is not the shared runtime alone, but the common user workflow it enables. In both use cases, the operator follows the same overall process: choose the use case, provide the required spatial and non-spatial inputs, confirm the interpreted result, execute it, and save or load the authored program.

That common path is why the thesis did not produce only two isolated demonstrations. It produced a prototype in which different task-oriented procedures can share the same acquisition, confirmation, execution, and persistence flow. The same structure also allows some components to change without forcing a redesign of the whole workflow, as long as the surrounding interfaces remain unchanged, although the current use case interfaces still limit how far that flexibility extends.

That result is still intentionally bounded. The system does not expose a full robot programming language, and it does not try to cover rich branching, loops, exception recovery, or complex sensor-driven logic. Its stronger result is narrower: some frequently changing tasks can be moved from repeated low-level program editing to re-authoring inside a use case that already defines most of the task structure.

\section{Pick-and-Place, Welding, and Use Case Scope}

The two demonstrated task families are not equivalent examples of the same success. They support different conclusions about where the approach is useful and what kind of task structure it fits best.

Pick-and-place is the clearest evidence that the prototype can combine authored intent with refreshed scene information at execution time. The authored command does not store only a fixed world-frame pick pose. It stores an object identity and a pick pose relative to that object, then reconstructs the world pick pose from a fresh scene estimate when the command is executed. This shows that the system is not limited to replaying fixed authored motions. This result is meaningful, but only within the limits of the current scene representation and the specific pick-and-place task demonstrated in the prototype.

Seam welding is the clearest practical fit for the present prototype. The authored trajectory expresses the main task intent directly in the workspace, and the use case can infer the surrounding robot procedure from that trajectory together with the detected scene geometry. In the demonstrated workflow, the user scans the scene, draws the weld, confirms the interpreted result, and then executes an approach--weld--depart sequence. That is a good match for the thesis question because it removes much of the low-level motion specification while keeping the authored input concrete and easy to interpret.

The usefulness of this workflow depends on how often the weld changes. For stable repetitive production, existing programming workflows may already be sufficient once the task has already been programmed. The present approach becomes more useful when the overall task stays similar but the concrete weld changes often enough that repeated re-teaching becomes a noticeable burden, for example in small-batch production or frequently reworked jobs. In that setting, the human provides the task-specific intent in context and the robot performs the repeated physical execution.

More broadly, the thesis supports selected use cases with explicit task definitions, not general robot programming. In those cases, most of the task structure is prepared in advance by developing the use case module, and the operator defines the concrete spatial parameters of each task instance in the real workspace. A narrow use case can expose only the spatial inputs that actually vary and keep authoring direct. A broader use case can cover more variation, but it must also expose more parameters and more internal choices. The thesis does not resolve the best balance between those two directions, but it does show that both can be built on the same runtime structure as long as the task remains inside a prepared use case rather than a full programming environment.

A workflow built around use cases could in principle be implemented without direct spatial input, for example by entering positions through a conventional interface. That would still simplify the developed system, but it would leave the main spatial part of robot authoring in a more indirect form. For the task families considered here, poses and trajectories are not secondary parameters. They are a central part of what changes between task instances, so direct spatial input is part of the simplification rather than an optional addition.

\section{Applicability Boundary and Limitations}

The first major boundary is scene detection. The current prototype detects only predefined, tagged, box-shaped objects and builds the scene description from that narrow representation. This was sufficient for the feasibility-oriented tests used in the thesis, but it is also the main unresolved question for any broader practical use of the system. The thesis does not establish how reliable or general this scene detection would be outside the demonstrated setup. That limitation matters directly for both representative tasks, because the practical usefulness of the overall workflow depends strongly on the quality of the scene information it can obtain.

The next boundary is sensing and calibration robustness. Pen tracking depends on marker visibility, lighting conditions, motion blur, and corner detection quality, and the current pen implementation does not fuse the available IMU data. Camera calibration also assumes that the stereo rig remains rigid with respect to the robot flange. These are workable assumptions for the thesis prototype, but they limit how confidently the present sensing stack can be transferred to less controlled conditions.

Execution remains equally bounded. The thesis demonstrates a shared robot motion interface and shows that use cases can drive real robot execution through it, but generalized tool control is not implemented. In the evaluated use cases, gripper and welding actions are only approximated with fixed delays. No motion planning, collision checking, or simulated execution is provided.

The current use cases also depend on task-specific operating assumptions. The weld command assumes that the workpiece configuration stays unchanged after confirmation. The pick-and-place command assumes that the scene can be refreshed at execution time and that the expected object can still be matched by its stored identity.

UI maturity and portability are secondary limits. The frontend prototype supports setup, authoring, execution, and persistence in one workflow, but it still shows clear prototype-level behavior in practice. Visualization can require manual reload, saved setups may restore only partially if hardware modules cannot reactivate, and hardware support remains platform-limited: camera acquisition is currently implemented only on Windows, while BLE discovery is currently functional on Windows but not on Linux.

At the same time, the developed system's modular structure provides a practical path for further development, because individual parts of the system can be improved without requiring the whole prototype to be rebuilt from the ground up. This does not remove the current limitations, and some improvements would still require changes at the interface boundary, but it does mean that the existing system remains a useful foundation for further work.

Taken together, these limits define the applicability boundary of the work. The thesis supports selected task-oriented workflows built around prepared use cases, not a general authoring method for arbitrary robot tasks. It does not yet provide the scene understanding, execution safety, or tooling maturity required for broader deployment claims.

The evaluation boundary is equally explicit. The thesis demonstrates qualitative feasibility with limited physical validation, and most development and testing took place in the Kassow REMII offline environment. It does not include a formal user study, a controlled comparison against teach pendant or offline programming workflows, or quantitative characterization of tracking accuracy, repeatability, or productivity. The result should therefore be read as a proof-of-concept for selected workflows, not as evidence of measured superiority or industrial readiness.
